\chapter{Alternative frontier: Sleep이 없어도 가능한 최신 경로}
\label{STC-S06}

Sleep-time compute을 promising하다고 판단하려면 이미 강한 방법을 약하게 묘사하면 안 된다. 이 장은
long context, retrieval, model editing, PEFT, context compilation, self-adaptation, recurrent memory
caching을 independent alternative 또는 component로 평가한다.

\section{Long context와 retrieval}

Long context는 원문을 손실 없이 한 번에 읽어야 하는 task의 default다. Context window와 attention
효율이 개선되면 sleep compilation의 일부 필요가 줄어든다. 그러나 매 query마다 같은 long history를
읽으면 token·KV·prefill cost가 반복되고, context selection과 stale conflict가 남는다. Sleep은 repeated
workload에서 hierarchy/summary/index를 만드는 build-time complement다.

Retrieval-Augmented Generation은 mutable knowledge를 weights 밖에 두고 retrieved passage로 generation을
condition하므로 update cost와 governance의 강한 baseline이다.\claim{STC-C026}
\parencite{srcstc0031} Record-level source, time, delete가 중요한 fact에서는 parametric sleep이 retrieval을
이기기 전에 정확도뿐 아니라 rollback과 rebase cost까지 이겨야 한다.

\section{Recurrent/neural memory와 serving cache}

RNN-style 또는 neural-memory model은 fixed-size state로 long history를 압축해 attention의 context cost를
낮출 수 있다. 그러나 prefix KV cache처럼 많은 request가 같은 reusable state를 공유하려면 state를
materialize하고 key를 정하며 cache/evict/share해야 한다. Memory Caching은 이 serving 문제를 명시한다.
\claim{STC-C025} \parencite{srcstc0022}

이는 sleep과 보완적이다. Wake model이 만드는 recurrent state를 session boundary에서 저장하고, sleep이
state를 compact/merge/rekey할 수 있다. 반대로 state reuse가 낮고 매 session이 unique하면 cache tier가
추가 complexity만 만든다. 평가에는 cache hit, state bytes, recompute cost, compatibility hash가 필요하다.

\section{PEFT: 작은 parametric destination}

LoRA는 adaptation을 low-rank weight delta로 제한해 trainable parameter와 optimizer-state cost를 줄이지만,
sequential interference와 lifecycle governance를 없애지는 않는다.\claim{STC-C027}
\parencite{srcstc0030,srcstc0036} User별 adapter가 작아도 수억 개가 되면 metadata lookup, base-version
compatibility, batching fragmentation, merge/delete가 새 serving problem이다.

PEFT는 broad full-weight sleep보다 유망한 parametric destination이다. Base를 immutable하게 두고 candidate
adapter를 version/disable하기 쉬우며, high-reuse behavior를 external profile보다 적은 wake overhead로
serve할 수 있다. 하지만 stable fact를 adapter에 넣으면 provenance와 temporal update가 어려워진다.

\section{Model editing: 좁은 update}

ROME은 transformer MLP의 factual association을 localized edit하여 targeted parametric change가 가능함을
보였지만 stable unbounded continual accumulation을 입증한 것은 아니다.\claim{STC-C028}
\parencite{srcstc0032} MEMIT은 많은 memories에 대한 mass editing을 확장했지만, 여전히 bounded batch
intervention이지 lifelong sleep loop의 증거는 아니다.\claim{STC-C029} \parencite{srcstc0033}

Editing은 correction latency가 중요하고 edit set이 작을 때 strong baseline이다. Sleep replay는 여러
episode에서 stable pattern을 통합할 때 strong하다. 양쪽 모두 efficacy 외에 paraphrase, locality, old edit
retention, multi-hop composition, rollback을 측정해야 한다.

\section{Context compilation과 one-forward adaptation}

Generative Adapter는 context를 한 forward pass로 parameter-like adapter에 compile하는 latent-memory 경로를
보인다. 그 가치는 이후 reuse가 compilation과 version management cost를 상각할 때 생긴다.
\claim{STC-C030} \parencite{srcstc0027} Gradient-based sleep보다 빠를 수 있으나 teacher/compiler quality,
base-version drift, adapter capacity, delete lineage가 남는다.

이 접근은 long context와 parametric memory의 중간점이다. 원문을 build path에서 한 번 읽고 compact
state로 serve한다. One-off document에는 compile하지 않고, predicted query count가 break-even을 넘는
document/user cluster만 선택하는 policy가 필요하다.

\section{Self-adaptation과 synthetic-data 위험}

Self-Adapting Language Models는 model이 adaptation data와 update rule을 스스로 생성해 wake learning과
later consolidation을 잇지만 recursive data quality risk를 만든다.\claim{STC-C031}
\parencite{srcstc0029,srcstc0041} Self-generated curriculum이 novelty를 늘리는지, 같은 error를 증폭하는지,
real evidence와 얼마나 연결되는지 cycle depth별로 평가한다.

\begin{table}[htbp]
\centering
\caption{Alternative frontier의 cost·capacity·governance 비교}
\label{tab:alternative-frontier}
\scriptsize
\begin{tabularx}{\textwidth}{p{0.14\textwidth} p{0.15\textwidth} p{0.15\textwidth} p{0.17\textwidth} X}
\toprule
방법 & One-time cost & Per-query cost & Capacity behavior & 주요 failure \\
\midrule
long context & 낮음 & token/KV/prefill 반복 & context window와 selection & distractor, latency \\
RAG/external & ingest/index & retrieval+context & independent storage growth & miss, poisoning, stale index \\
recurrent state & state build & small state forward & fixed/compressed state & state collision, low reuse \\
LoRA/adapter & training+validation & routing+delta compute & artifact bank 증가 & interference, fragmentation \\
model editing & locate/edit/test & 거의 없음 & edits 누적 & locality, composition \\
context compile & compiler inference & adapter reuse & compiled artifact 증가 & break-even, rebase \\
global refresh & large periodic train & model-only inference & release 단위 재압축 & freshness lag, high cost \\
\bottomrule
\end{tabularx}
\end{table}

\begin{keypoint}[Alternative와 STC의 관계]
Sleep은 별도 optimizer family가 아니다. Retrieval index rebuild, recurrent-state compaction, adapter
training, model editing, distillation, global refresh를 언제·어떤 data horizon에서 실행하고 어떻게
publish할지 묶는 lifecycle다. 따라서 ``STC vs RAG''보다 ``RAG-only vs background synthesis vs selective
promotion'' 같은 system configuration을 비교해야 한다.
\end{keypoint}
